\chapter{Background, aim and objectives}

\section{Significance of the project}

\instructionbox{Explain, for a reader outside your group, why this project
is worth doing. Set the scene: what is the broader scientific, technical or
societal context? What existing gap, problem or open question motivates
your work? Aim for roughly half a page to one page. Delete this box and
replace the placeholder text below with your own content.}

Describe the wider context and motivation for the project here. What field
does the project belong to, what is currently known, and why does it
matter? Keep this section understandable to a reader who is not a
specialist in the topic.

\newpage

\section{Research questions and/or goals}

\instructionbox{State the concrete questions your project sets out to
answer, or the concrete goals it aims to reach. Bullet points work well
here. Each question/goal should be specific enough that you can later
judge, objectively, whether it was answered or reached. Avoid vague goals
such as ``learn more about X'' -- prefer ``determine how X affects Y under
conditions Z''.}

The main goal of this project is to \emph{[state the overall aim in one
sentence]}. Our specific research questions are:

\begin{itemize}
    \item \emph{[Research question 1]}
    \item \emph{[Research question 2]}
    \item \emph{[Research question 3]}
\end{itemize}

\section{Expected results and their potential impact}

\instructionbox{Describe what you expect to find out or produce, and who
would benefit from it (your field, future students, an external
collaborator, society more broadly). This section should connect back to
the significance described above -- i.e.\ explain how reaching your goals
addresses the motivation you gave earlier.}

Describe the expected outcomes of the project and their relevance here.

